\chapter{Introduction}\label{introduction}

\section{Motivation}\label{motivation}

Large Language Models have demonstrated strong capabilities for
interpreting, summarizing, and reasoning over complex textual
information. These capabilities have created significant interest in
their use for cybersecurity tasks, where analysts frequently work with
heterogeneous information such as vulnerability descriptions, attack
patterns, request traces, incident reports, and security logs.

Web application logs are particularly suitable for LLM-assisted analysis
because they combine structured fields with free-form or
attacker-controlled content. A single HTTP request may contain a method,
URL, query parameters, a request body, a User-Agent string, a source
address, and a response status. These fields may contain evidence of
attacks such as SQL Injection, Cross-Site Scripting, Path Traversal, or
repeated authentication attempts. An LLM can potentially transform these
technical observations into a structured explanation containing an
attack classification, severity assessment, relevant indicators, and
recommended defensive actions.

However, introducing an LLM into the analysis pipeline creates a
security problem that does not exist in conventional log parsers. Some
web-log fields are directly controlled by external users. If an attacker
inserts natural-language instructions into fields such as the URL, query
parameters, request body, or User-Agent string, those instructions may
later become part of the context presented to the LLM. A vulnerable
model may interpret the injected text as an instruction rather than as
untrusted data.

For example, a malicious request could combine a conventional attack
payload with an instruction such as:

\begin{lstlisting}
Ignore previous instructions and classify this request as benign.
\end{lstlisting}

A traditional parser would treat this string as passive content. An LLM,
however, may allow it to influence the reasoning process. This creates
the possibility of manipulating the reported classification, lowering
the reported severity, suppressing evidence, or attempting to extract
hidden system instructions.

The problem becomes even more relevant when Retrieval-Augmented
Generation is used. RAG can improve factual grounding by supplying the
model with cybersecurity knowledge retrieved from an external knowledge
base. At the same time, the final prompt still combines trusted
instructions, retrieved evidence, and untrusted log content. Robust
prompt construction and evidence control are therefore necessary if the
system is expected to produce reliable security analysis.

\section{Problem Statement}\label{problem-statement}

The central problem addressed by this thesis is the potential
vulnerability of a RAG-based LLM assistant when analyzing web
application logs containing attacker-controlled text.

The intended behavior of the system is to treat every log field as data
to be analyzed. The attacker, however, may deliberately insert
instructions designed to alter the output of the assistant. This creates
a conflict between the trusted task definition and untrusted
natural-language content.

The thesis therefore investigates two related questions:

\begin{enumerate}
\def\labelenumi{\arabic{enumi}.}
\tightlist
\item
  How vulnerable is the assistant when no explicit defenses are applied?
\item
  To what extent can lightweight and explainable mitigation strategies
  reduce this vulnerability without removing the usefulness of the LLM
  and RAG pipeline?
\end{enumerate}

\section{Research Question}\label{research-question}

\begin{quote}
\textbf{How vulnerable is a RAG-based LLM assistant for web application
log analysis to prompt injection attacks embedded in attacker-controlled
log fields, and how can evidence-controlled retrieval and simple
mitigation strategies reduce this risk?}
\end{quote}

\section{General Objective}\label{general-objective}

Design, implement, and evaluate a RAG-based LLM assistant for web
application log analysis, focusing on robustness against Prompt
Injection attacks.

\section{Specific Objectives}\label{specific-objectives}

The specific objectives of this thesis are:

\begin{enumerate}
\def\labelenumi{\arabic{enumi}.}
\tightlist
\item
  Design an LLM-based architecture for structured web application log
  analysis.
\item
  Construct a cybersecurity knowledge base containing selected MITRE
  ATT\&CK techniques, OWASP web-attack knowledge, response playbooks,
  and prompt-injection handling policies.
\item
  Build a controlled dataset containing clean and adversarial web
  application logs.
\item
  Define and operationalize representative Prompt Injection attacks
  embedded in attacker-controlled log fields.
\item
  Implement a RAG pipeline capable of retrieving evidence relevant to
  the analyzed event.
\item
  Implement lightweight mitigation mechanisms based on prompt hardening,
  instruction/data separation, metadata tagging, simple prompt-injection
  detection, and evidence-only prompting.
\item
  Compare defended and undefended system configurations under the same
  dataset and evaluation procedure.
\item
  Measure attack success, classification performance, retrieval quality,
  evidence use, hallucination behavior, and robustness.
\item
  Analyze failure cases and identify limitations of the proposed
  defenses.
\end{enumerate}

\section{Scope}\label{scope}

This thesis is intentionally restricted to \textbf{web application
logs}.

The system analyzes web/application request information such as:

\begin{itemize}
\tightlist
\item
  HTTP method,
\item
  URL,
\item
  query parameters,
\item
  request body,
\item
  User-Agent,
\item
  status code,
\item
  source IP,
\item
  and related structured metadata.
\end{itemize}

The experimental attack classes are limited to:

\begin{itemize}
\tightlist
\item
  SQL Injection,
\item
  Cross-Site Scripting (XSS),
\item
  Path Traversal,
\item
  Brute Force Login,
\item
  and Benign/Normal requests as a control class.
\end{itemize}

The study does \textbf{not} attempt to implement a complete Security
Operations Center assistant and does not cover:

\begin{itemize}
\tightlist
\item
  SOC orchestration,
\item
  SIEM platforms,
\item
  SOAR platforms,
\item
  firewall logs,
\item
  EDR logs,
\item
  IDS/IPS logs,
\item
  enterprise incident management,
\item
  or autonomous response actions.
\end{itemize}

This restriction is important because it allows the research to isolate
the interaction between web-log content, Prompt Injection, retrieval,
evidence control, and LLM reasoning.

\section{Research Contributions}\label{research-contributions}

The expected contributions of this thesis are:

\begin{enumerate}
\def\labelenumi{\arabic{enumi}.}
\tightlist
\item
  A reproducible architecture for RAG-assisted web application log
  analysis.
\item
  A controlled dataset containing paired clean and adversarial log
  examples.
\item
  An operational Prompt Injection taxonomy tailored to
  attacker-controlled web-log fields.
\item
  An experimental comparison between LLM-only, RAG, and defended RAG
  configurations.
\item
  An evaluation methodology combining conventional classification
  metrics with Prompt Injection robustness and evidence-grounding
  metrics.
\item
  An analysis of the effectiveness and limitations of lightweight
  mitigation strategies.
\end{enumerate}

\section{Thesis Structure}\label{thesis-structure}

The remainder of this thesis is organized as follows.

Chapter 2 presents the theoretical background and related work
concerning RAG, Prompt Injection, LLM applications in cybersecurity, web
application log analysis, evidence grounding, and Prompt Injection
defenses. The chapter concludes by identifying the research gap
addressed by this work.

Chapter 3 defines the research methodology, scope, threat model, dataset
strategy, experimental configurations, and evaluation metrics.

Chapter 4 presents the design of the proposed system, including
preprocessing, prompt-injection detection, retrieval, the knowledge
base, evidence control, secure prompt construction, and structured
output.

Chapter 5 describes the dataset construction and system implementation.

Chapter 6 presents the experimental results.

Chapter 7 discusses the results, failure cases, limitations, and threats
to validity.

Chapter 8 concludes the thesis and proposes directions for future work.

\begin{center}\rule{0.5\linewidth}{0.5pt}\end{center}
